\documentclass[10pt,a4paper]{amsart}
\usepackage[margin=19mm]{geometry}
\usepackage{amsmath,amssymb,mathtools}
\usepackage[scr=rsfs,cal=euler]{mathalfa}
\usepackage{xfrac}
\usepackage[unicode,hidelinks,bookmarks=false]{hyperref}
\newcommand{\ie}{i.e.\ }
\newcommand{\cf}{cf.\ }
\newcommand{\resp}{respectively\ }
\newcommand{\Subsection}{\subsection}
\providecommand{\nobreakdash}{\nobreak\hbox{-}}
\setlength{\emergencystretch}{2em}
\multlinegap0pt
\usepackage{bm}
\usepackage{bbm}
\usepackage{dsfont}
\usepackage{xspace}
\usepackage{cancel}
\usepackage{mathtools}
\usepackage{amsthm}
\makeatletter
\@ifundefined{theorem}{\newtheorem{theorem}{Theorem}}{}
\@ifundefined{lem}{\newtheorem{lem}{Lemma}}{}
\@ifundefined{proposition}{\newtheorem{proposition}[theorem]{Proposition}}{}
\makeatother
\DeclareMathOperator{\rank}{rank}
\newcommand{\bbQ}{\mathbb{Q}}
\newcommand{\legendre}[2]{\genfrac{(}{)}{}{}{#1}{#2}}
\title[Classification of the rank of a certain family of elliptic c]{Classification of the rank of a certain family of elliptic curves}
\author{Arkabrata Ghosh, Bidisha Roy,, Richa Sharma}
\date{Source: arXiv:2607.16335v2 (CC BY 4.0)}
\begin{document}
\maketitle

\noindent\emph{Source and licence.} Classification of the rank of a certain family of elliptic curves. Version arXiv:2607.16335v2 (CC BY 4.0), distributed under the Creative Commons Attribution 4.0 International licence, as recorded in the arXiv licence metadata for this exact version and shown on the abstract page https://arxiv.org/abs/2607.16335v2. Full rights notice: licenses/arxiv-2607.16335-CC-BY-4.0.txt.\par\smallskip

\noindent\textit{Editorial scope.} Counted unit: the Lem whose source label is \texttt{lem:oalpha\_3.1} in the article (source0.tex lines 1419--1431, source label \texttt{lem:oalpha\_3.1}), delivered together with the article's own complete original proof of it (source0.tex lines 1836--1876, one \texttt{proof} environment of 41 lines). No mathematics is written, repaired or re-derived: statement and proof are the article's own text line for line. The proof is detached from the statement in the source: the article states the result at lines 1419--1431 and prints its proof at lines 1836--1876, and the proof environment's own header names the statement, so the association is the article's own. Required context retained uncounted: prop:2.1 (lines 370--378); lem:root (lines 392--398); lem:alpha\_case\_1.2 (lines 751--763); lem:alpha\_case\_1.3 (lines 771--783); lem:alpha\_case\_2.1 (lines 816--828); lem:alpha\_case\_2.2 (lines 856--868); lem:alpha\_case\_3.1 (lines 932--944); lem:alpha\_case\_3.3 (lines 963--975); lem:oalpha\_1.2 (lines 1256--1268); lem:oalpha\_1.3 (lines 1280--1292); lem:oalpha\_2.1 (lines 1325--1337); lem:oalpha\_2.2 (lines 1350--1362); lem:oalpha\_3.3 (lines 1457--1469). Every definition, display and label of those lines is retained unchanged. Omitted from this package and not redistributed: 1--369, 379--391, 399--750, 764--770, 784--815, 829--855, 869--931, 945--962, 976--1255, 1269--1279, 1293--1324, 1338--1349, 1363--1418, 1432--1456, 1470--1835, 1877--2725. Nothing inside a retained excerpt is omitted or altered: every retained body is delivered exactly as the article's lines are written, so no omission receipt of any kind accompanies this package. Citations: the retained excerpts carry no citation command at all, so no bibliography entry of the article is transcribed and no reference is redistributed. Reference adaptation: none was needed and none was made; every reference inside the delivered unit resolves inside it, so no unresolved reference or citation is produced. Printed slips kept literal: no printed word, symbol, hypothesis or number of the counted statement or of its proof was altered. Layout and class: the article's own class is \texttt{amsart} with packages including amssymb, xcolor, fullpage, hyperref, listings, ulem, minted, comment; the wrapper is the lane's pinned \texttt{amsart} editorial wrapper at 10pt on A4, which restates the theorem-like environments the retained text needs and adds the article's own macro definitions that the retained text needs (\texttt{\textbackslash bbQ}, \texttt{\textbackslash legendre}, \texttt{\textbackslash rank}), with extra wrapper packages (bm, bbm, dsfont, xspace, cancel, mathtools). The delivered numbering is the unit's own. Assets: no retained span contains a figure environment, a graphics-inclusion command, a PDF-page inclusion, a table or a TikZ picture; no image file is copied into this package, the package declares no asset dependency and no image file is redistributed with it. Licence: version arXiv:2607.16335v2 of this article is distributed under the Creative Commons Attribution 4.0 International licence, as recorded in the arXiv licence metadata for this exact version and shown on the abstract page https://arxiv.org/abs/2607.16335v2. A full rights notice is delivered at licenses/arxiv-2607.16335-CC-BY-4.0.txt. Characters: every retained excerpt is pure ASCII, so the delivered engine, XeLaTeX with its default Latin Modern text font, reports no missing character.
\begin{proposition}
\label{prop:2.1}
Let $r$ be the rank of $E(\bbQ)$, where $E$, $\alpha$ and $\overline{\alpha}$
are as above. Then,
\[
\frac{1}{4} |\alpha(E(\bbQ))| \cdot \left| \overline{\alpha} \left( \overline{E}(\bbQ) \right) \right|
= 2^{r}.
\]
\end{proposition}

\begin{lem}
\label{lem:root}
Let $p$ and $q$ be distinct odd primes. Then
\[
w\left(E_{-2pq}/\bbQ\right)=-1.
\]
\end{lem}

\begin{lem}
\label{lem:alpha_case_1.2}
If $(p,q)\equiv (1,7)\pmod 8$, then
\[
\alpha\left(E_{-2pq}(\bbQ)\right)\subseteq
\begin{cases}
\{1,2,p,-q,2p,-2q,-pq,-2pq\},
& \text{for } \legendre{p}{q}=1,\\
\{1,2,-pq,-2pq\},
& \text{for } \legendre{p}{q}=-1.
\end{cases}
\]
\end{lem}

\begin{lem}
\label{lem:alpha_case_1.3}
If $(p,q)\equiv (7,1)\pmod 8$, then
\[
\alpha\left(E_{-2pq}(\bbQ)\right)\subseteq
\begin{cases}
\{1,2,-p,q,-2p,2q,-pq,-2pq\},
& \text{for } \legendre{p}{q}=1,\\
\{1,2,-pq,-2pq\},
& \text{for } \legendre{p}{q}=-1.
\end{cases}
\]
\end{lem}

\begin{lem}
\label{lem:alpha_case_2.1}
If $(p,q)\equiv (1,3)\pmod 8$, then
\[
\alpha\left(E_{-2pq}(\bbQ)\right)\subseteq
\begin{cases}
\{1,-2,p,q,-2p,-2q,pq,-2pq\},
& \text{for } \legendre{p}{q}=1,\\
\{1,-2,pq,-2pq\},
& \text{for } \legendre{p}{q}=-1.
\end{cases}
\]
\end{lem}

\begin{lem}
\label{lem:alpha_case_2.2}
If $(p,q)\equiv (1,5)\pmod 8$, then
\[
\alpha\left(E_{-2pq}(\bbQ)\right)\subseteq
\begin{cases}
\{1,-1,p,-p,2q,-2q,2pq,-2pq\},
& \text{for } \legendre{p}{q}=1,\\
\{1,-1,2pq,-2pq\},
& \text{for } \legendre{p}{q}=-1.
\end{cases}
\]
\end{lem}

\begin{lem}
\label{lem:alpha_case_3.1}
If $(p,q)\equiv (3,1)\pmod 8$, then
\[
\alpha\left(E_{-2pq}(\bbQ)\right)\subseteq
\begin{cases}
\{1,-2,p,q,-2p,-2q,pq,-2pq\},
& \text{for } \legendre{p}{q}=1,\\
\{1,-2,pq,-2pq\},
& \text{for } \legendre{p}{q}=-1.
\end{cases}
\]
\end{lem}

\begin{lem}
\label{lem:alpha_case_3.3}
If $(p,q)\equiv (5,1)\pmod 8$, then
\[
\alpha\left(E_{-2pq}(\bbQ)\right)\subseteq
\begin{cases}
\{1,-1,q,-q,2p,-2p,2pq,-2pq\},
& \text{for } \legendre{p}{q}=1,\\
\{1,-1,2pq,-2pq\},
& \text{for } \legendre{p}{q}=-1.
\end{cases}
\]
\end{lem}

\begin{lem}
\label{lem:oalpha_1.2} 
If $(p,q)\equiv (1,7)\pmod 8$, then
\[
\overline{\alpha}\left(\overline{E_{-2pq}}(\bbQ)\right)
\begin{aligned}
    &\subseteq \{1,p,2q,2pq\},
&& \text{for } \legendre{p}{q}=1,\\
&= \{1,2pq\},
&& \text{for } \legendre{p}{q}=-1.
\end{aligned}
\]
\end{lem}

\begin{lem}
\label{lem:oalpha_1.3} 
If $(p,q)\equiv (7,1)\pmod 8$, then
\[
\overline{\alpha}\left(\overline{E_{-2pq}}(\bbQ)\right)
\begin{aligned}
&\subseteq \{1,q,2p,2pq\},
&& \text{for } \legendre{p}{q}=1,\\
&= \{1,2pq\},
&& \text{for } \legendre{p}{q}=-1.
\end{aligned}
\]
\end{lem}

\begin{lem}
\label{lem:oalpha_2.1}
If $(p,q)\equiv (1,3)\pmod 8$, then
\[
\overline{\alpha}\left(\overline{E_{-2pq}}(\bbQ)\right)
\begin{aligned}
&\subseteq \{1,p,2q,2pq\},
&& \text{for } \legendre{p}{q}=1,\\
&= \{1,2pq\},
&& \text{for } \legendre{p}{q}=-1.
\end{aligned}
\]
\end{lem}

\begin{lem}
\label{lem:oalpha_2.2}
If $(p,q)\equiv (1,5)\pmod 8$, then
\[
\overline{\alpha}\left(\overline{E_{-2pq}}(\bbQ)\right)
\begin{aligned}
&\subseteq \{1,p,2q,2pq\},
&& \text{for } \legendre{p}{q}=1,\\
&= \{1,2pq\},
&& \text{for } \legendre{p}{q}=-1.
\end{aligned}
\]
\end{lem}

\begin{lem}
\label{lem:oalpha_3.3}
If $(p,q)\equiv (5,1)\pmod 8$, then
\[
\overline{\alpha}\left(\overline{E_{-2pq}}(\bbQ)\right)
\begin{aligned}
&\subseteq \{1,q,2p,2pq\},
&& \text{for } \legendre{p}{q}=1,\\
&= \{1,2pq\},
&& \text{for } \legendre{p}{q}=-1.
\end{aligned}
\]
\end{lem}

\begin{lem}
\label{lem:oalpha_3.1}
If $(p,q)\equiv (3,1)\pmod 8$, then
\[
\overline{\alpha}\left(\overline{E_{-2pq}}(\bbQ)\right)
\begin{aligned}
&\subseteq \{1,q,2p,2pq\},
&& \text{for } \legendre{p}{q}=1,\\
&= \{1,2pq\},
&& \text{for } \legendre{p}{q}=-1.
\end{aligned}
\]
\end{lem}

\begin{proof}
For the six residue classes occurring in the theorem, the relevant pairs of
lemmas are
\[
\begin{array}{c|c}
(p,q)\pmod8 & \text{lemmas used}\\
\hline
(1,3) & \ref{lem:alpha_case_2.1},\ \ref{lem:oalpha_2.1}\\
(3,1) & \ref{lem:alpha_case_3.1},\ \ref{lem:oalpha_3.1}\\
(1,5) & \ref{lem:alpha_case_2.2},\ \ref{lem:oalpha_2.2}\\
(5,1) & \ref{lem:alpha_case_3.3},\ \ref{lem:oalpha_3.3}\\
(1,7) & \ref{lem:alpha_case_1.2},\ \ref{lem:oalpha_1.2}\\
(7,1) & \ref{lem:alpha_case_1.3},\ \ref{lem:oalpha_1.3}
\end{array}
\]
In every row, the corresponding lemmas give
\[
\begin{aligned}
|\alpha(E_{-2pq}(\bbQ))|&\leq8,
&\left|\overline{\alpha}\left(\overline{E_{-2pq}}(\bbQ)\right)\right|&\leq4,
&&\text{for }\legendre{p}{q}=1,\\
|\alpha(E_{-2pq}(\bbQ))|&\leq4,
&\left|\overline{\alpha}\left(\overline{E_{-2pq}}(\bbQ)\right)\right|&=2,
&&\text{for }\legendre{p}{q}=-1.
\end{aligned}
\]
Thus Proposition~\ref{prop:2.1} yields
$\displaystyle
r\leq3 \quad\text{for }\legendre{p}{q}=1,
\qquad
r\leq1 \quad\text{for }\legendre{p}{q}=-1.
$
By Lemma~\ref{lem:root} and the Parity Conjecture, $r$ is odd. Hence
\[
\rank(E_{-2pq}(\bbQ))=
\begin{cases}
1\text{ or }3, & \text{for }\legendre{p}{q}=1,\\
1, & \text{for }\legendre{p}{q}=-1.
\end{cases}
\]
\end{proof}

\end{document}
